%%%%%%%%%%%%%%%%%%%%%%%%%%%%%%%%%%%%%%%%%%%%%%%%%%%%%%%%%%
% # Introduction
%%%%%%%%%%%%%%%%%%%%%%%%%%%%%%%%%%%%%%%%%%%%%%%%%%%%%%%%%%

\addcontentsline{toc}{part}{Chapters}
\chapter{Introduction}\label{ch-introduction}

Second language (L2) acquisition is a complex and fascinating process. In addition to facing many of the same challenges that children do when acquiring a first language (L1), L2 learners (L2ers) must contend with the added complication of having to integrate information about the new language into a broader cognitive system that already has at least one other language firmly rooted in it. L2 acquisition also differs from L1 acquisition in the amount and type of input that learners receive as well as in the maturational state of the brain at the time of acquisition. In light of differences such as these, the difficulty that adult L2ers typically experience is perfectly understandable, and the fact that so many achieve high levels of communicative competence in their nonnative languages is a truly impressive feat. 

This dissertation investigates several important theoretical issues in the L1 and L2 linguistics literature by looking at the processing and acceptability of relative clauses (RCs) among English native speakers, L1-Korean L2ers of English, and L1-Mandarin L2ers of English. Specifically, I compare gapped RCs like the one in (\ref{ex-int-rel-a}) to resumptive RCs like the one in (\ref{ex-int-rel-b}).

\pex
\a the man that Tina hired\label{ex-int-rel-a}
\a \ljudge*the man that Tina hired him\label{ex-int-rel-b}
\xe 

\noindent Whereas English and Korean are commonly thought to prohibit resumption in all the RC environments I examine (for~English, see \cite{beltrama.xiang2016}; \cite{han.etal2012}; \cite{keffala.goodall2011}; \cite{mcdaniel.cowart1999}; \cite{omaki.nakao2010}; \cite{polinsky.etal2013}; for~Korean, see \cite{han2013}; \cite{keenan.comrie1977}; \cite{kim2013}; \cite{song2003}), Mandarin is thought to allow resumption in all but one of those same environments (\cite{keenan.comrie1977}; \cite{li.wei2014}; \cite{pan2016}).

What makes resumptive RCs a topic of special interest among L2 acquisition researchers is that it belongs to a family of relatively rare morphosyntactic phenomena that are known to crop up from time to time in L2 spoken and written output even when they are disallowed in both the L1 and the target language. In such cases, the emergence of the phenomenon cannot be straightforwardly traced to either target-language input or L1 transfer, and so it must have a source internal to the learner (\cite{eckman2004}, p.~89; \cite{schulz2006}, p.~1). Studies on resumption in RCs are thus ideal for the investigation of theoretical issues at the intersection of linguistic competence and linguistic performance (\cite{meltzer-asscher2021}, p.~191). 

Previous L2 experimental work has suggested that when resumptive RCs cannot be traced to L1 influence, they may instead arise because they facilitate processing when cognitive load is high (\cite{algady2013}; \cite{hyltenstam1984,hyltenstam1990}; \cite{maghrabi1997}); however, this has only been looked at in production and not in comprehension. It is often assumed that such resumptive RCs represent grammatical relativization in the interlanguage grammars of L2ers who produce them (\cite{eckman2004,hyltenstam1984,hyltenstam1990}), but this inference, too, needs to be tested. Existing L2 studies on resumption in (mostly English) RCs have been limited in a number of other ways: For instance, they have almost exclusively focused on short-distance RCs (i.e.,~within a single clause, as in (1)), and none has directly tested how resumption affects accurate RC interpretation.

The present dissertation addresses these gaps in the research by drawing on innovative methods and technologies to probe the production, comprehension, and acceptability of English gapped vs. resumptive RCs across three increasingly complex syntactic environments: short-distance, long-distance, and \textit{wh}-island. In doing so, it explores key questions related to L2 knowledge, L2 use, and the development of each.

%========================================================%
\section{Theoretical issues explored in this dissertation}
%========================================================%

This dissertation explores four main theoretical issues, most of which have implications that go beyond the narrow investigation of gapping and resumption in RCs.

%--------------------------------------------------------%
\subsection{RC processing in production vs. comprehension}
%--------------------------------------------------------%

While there is broad consensus that resumption assists with RC production under processing strain (e.g.,~\cite{asudeh2004}; \cite{ferreira.swets2005}; \cite{zukowski.larsen2004}), whether it can also help with RC comprehension is more controversial. There are a number of existing L1 adult English studies that claim to show that resumption makes RCs easier to understand (\cite{beltrama.xiang2016}; \cite{hammerly2022}; \cite{hofmeister.norcliffe2013}), but almost none has checked for accurate RC interpretation. Proponents of the selfish speaker model (e.g., \cite{ferreira.swets2005}; \cite{polinsky.etal2013}; \cite{morgan.etal2020}) argue that resumption actually hinders RC comprehension, taking this as a sign that there is a deep rift between the language production system and the language comprehension system. I put the selfish speaker model to the test in the present dissertation by examining the effect that resumption has on both RC production and RC comprehension.

%--------------------------------------------------------%
\subsection{The status of resumptive RCs in English interlanguages}
%--------------------------------------------------------%

Previous experimental studies have shown that L2ers of English often accept resumptive RCs at higher rates than native English speakers do; the contrast is especially pronounced in cases where the L1 has grammatical resumption in the relevant syntactic environment (\cite{alosaimi2021,hawkins.chan1997,maghrabi1997,marefat.abdollahnejad2014,solaimani.etal2023,tezel1998}), but it is also present in cases, albeit to a lesser extent, where the L1 lacks grammatical resumption in the relevant environment (\cite{hawkins.chan1997,tezel1998}). I use acceptability judgment tasks to investigate whether resumption represents a licit relativization option in the English interlanguages of my participants, interpreting consistent (≥~80\%) acceptance of RP trials as a sign that L2ers may have an interlanguage grammar that allows resumptive RCs in the relevant environment. For individuals who do consistently accept resumptive RCs, I probe for L1 transfer effects by collecting ratings from L2ers not only in English but also in their L1 (Korean or Mandarin). Checking that L2ers perform as expected in the L1 (based on what is reported in the literature as well as native-speaker judgments) is an important step to include in L2 research of this type, and yet it is one that has only very seldom been included in previous L2 studies of any kind (but~see \cite{zenker.schwartz2017}).

%--------------------------------------------------------%
\subsection{The relationship between competence and performance}
%--------------------------------------------------------%

It has traditionally been assumed in the L2 literature that L2ers' production patterns reflect representations licensed by their interlanguage grammars. While this assumption, which I call the Performance-Competence Correspondence Assumption (PCCA), has been questioned (e.g.,~\cite{kim2016}, p.~4ff; \cite{schwartz1986}, p.~122ff; \cite{schwartz1993}, p.~150), it is still implicitly if not explicitly assumed in many L2 studies, especially those on L2 gapped vs. resumptive RCs (\cite{algady2013}, p.~184; \cite{eckman2004}, p.~106ff; \cite{hyltenstam1984}, p.~30; \cite{marefat.abdollahnejad2014}, p.~36; \cite{pavesi1986}, p.~49). I test the PCCA by examining whether unambiguous processing facilitation effects in the production and comprehension of resumptive RCs are limited to participants who consistently rate them as acceptable in one or more of the RC environments tested.

%--------------------------------------------------------%
\subsection{The purging of nontarget patterns from interlanguages}
%--------------------------------------------------------%

If resumption represents a licit relativization option for at least some L2ers, then the question arises as to whether it can ever be purged from the interlanguage in cases where the target language does not permit it. Previous L2 English studies have shown that for L1-Arabic L2ers (\cite{alosaimi2021,maghrabi1997}), L1-Cantonese L2ers (\cite{hawkins.chan1997}), L1-French L2ers (\cite{hawkins.chan1997,solaimani.etal2023}), and L1-Persian L2ers (\cite{marefat.abdollahnejad2014,solaimani.etal2023}), higher proficiency scores are associated with lower rates of acceptance for resumptive RCs, suggesting that such purging is indeed possible. In this dissertation, I test whether a sensitivity to the English prohibition on resumption in RCs can develop over time in the relevant L2ers.

%========================================================%
\section{Organization of the dissertation}
%========================================================%

The remainder of the dissertation is organized as follows: The next section, which runs from Chapter~\ref{ch-properties} through Chapter~\ref{ch-acquisition}, lays the conceptual groundwork for this research. Chapter~\ref{ch-properties} provides general background information on gapped and resumptive RCs, paying special attention to my three languages of interest: English, Korean, and Mandarin. Chapter~\ref{ch-gapping} discusses some of the factors that are thought to influence the ease with which gapped RCs are processed, focusing on the three-way short-distance vs. long-distance vs. \textit{wh}-island contrast that I examine in my experiments. Chapter~\ref{ch-resumption} reviews the L1 adult literature on resumption in RCs. Chapter~\ref{ch-acquisition} does the same for research in the L2 context, followed by exploring how learners develop L1 vs. L2 implicit knowledge regarding the distribution of gapped vs. resumptive RCs. Along the way, I introduce the theoretical issues that motivated the present research.

The experimental portion of the dissertation runs from Chapter~\ref{ch-study} through Chapter~\ref{ch-reanalysis}. Chapter~\ref{ch-study} introduces the present study by describing the tasks, participants, and procedures in my experiments. The next three chapters report the primary results from each of the main tasks in turn: first the elicited production task (EPT; Chapter~\ref{ch-ept}), then the self-paced reading task (SPRT; Chapter~\ref{ch-spr}), and finally the acceptability judgment tasks (AJTs; Chapter~\ref{ch-ajt}). The last chapter in this part of the dissertation (Chapter~\ref{ch-reanalysis}) reexamines, with consistent acceptors of English resumptive RCs removed, the processing data from the EPT and the SPRT to test whether participants must consider resumptive RCs acceptable in order to derive a processing benefit from them.

Chapter~\ref{ch-discussion} is the final chapter in the dissertation. It begins by reviewing the major findings from the experiments and discussing their theoretical implications as they pertain to the central issues explored in the project. It then goes on to consider limitations of the present study and to make several recommendations regarding future avenues of research.